\section{Notes on the PDF}
\label{sec:notes}

An author can also ask for changes directly on the PDF: select text, or mark a
region around a figure or an equation, and write a note. All pending notes go
to the agent as one task. For that task to be useful, each note must name the
source it refers to.

\subsection{From a mark to source lines}

A note is a set of rectangles $\{r_1, \dots, r_q\}$ in the normalised
coordinates of \cref{eq:normalised} on page $P$, together with a comment. For
every printed line the selection touches, Papilot asks SyncTeX for the source
position of the first and the last marked character~\cite{laurens2008}:
\begin{equation}
  \mathcal{S}(P, x, y) = (\text{file}, \text{line}),
  \qquad
  L = \{\, \mathcal{S}(P, c) : c \in C_{\text{first}} \cup C_{\text{last}} \,\},
  \qquad
  [\ell_{-}, \ell_{+}] = [\min L, \max L].
  \label{eq:synctex-range}
\end{equation}

SyncTeX records boxes and glue rather than characters. The space before a word
that begins a printed line comes from the line break of the \emph{previous}
source line, so such a word is credited to that line. Papilot therefore
\emph{settles} each end of the selection by matching its printed context
against the source near SyncTeX's answer.

\subsection{Settling a selection}

Let $\letters(\cdot)$ map text to its lowercase letters and digits, after
Unicode compatibility normalisation, with \TeX{} control sequences and comments
removed from the source. For a source window $[\ell_{-} - 4, \ell_{+} + 4]$ we
form the string $S = \letters(\text{lines})$ and remember for each position the
line it came from, $\pi : \{1, \dots, \abs{S}\} \to \N$. The distance of a
line from SyncTeX's range is
\begin{equation}
  \dist(\ell) = \max\bigl(0,\ \ell_{-} - \ell,\ \ell - \ell_{+}\bigr).
  \label{eq:line-distance}
\end{equation}

\begin{algorithm}[t]
\caption{Settling one end of a selection}
\label{alg:settle}
\begin{algorithmic}[1]
\Require printed line $p$, index $\alpha$ of the marked letter in $\letters(p)$, source string $S$ with line map $\pi$
\For{$\rho \in (\infty, 24, 10)$} \Comment{whole line first, then shorter context}
  \State $o_0 \gets \max(0, \alpha - \rho)$
  \State $q \gets \letters(p)_{[o_0,\ \alpha + \rho]}$
  \State $O \gets \{ o : S_{[o,\, o + \abs{q})} = q \}$
  \If{$O \neq \emptyset$}
    \State $o^{\star} \gets \argmin_{o \in O} \dist\bigl(\pi(o + \alpha - o_0)\bigr)$
    \State \Return $\pi\bigl(o^{\star} + \alpha - o_0\bigr)$
  \EndIf
\EndFor
\State \Return SyncTeX's line \Comment{formulas and other unmatched text}
\end{algorithmic}
\end{algorithm}

\Cref{alg:settle} runs once for the first and once for the last marked
character. If the settled range is empty or inverted, Papilot keeps SyncTeX's
range. Because the comparison uses letters only, subscripts, macros and
hyphenation do not prevent a match; formulas, whose printed form shares no
letters with their source, fall back to SyncTeX by design.

\subsection{When is a match unique?}

\Cref{alg:settle} picks the occurrence nearest to SyncTeX's answer, so a
spurious occurrence matters only if it is closer than the true one. Model the
letters of the source as a stationary source with entropy $\mathcal{H}$ bits per letter.
The expected number of chance occurrences of a context of $\abs{q}$ letters in
a window of $\abs{S}$ letters is about
\begin{equation}
  \E[\text{spurious}] \approx \abs{S} \cdot 2^{-\mathcal{H} \abs{q}} .
  \label{eq:spurious}
\end{equation}
Estimates of the entropy of English range from about one to four bits per
letter depending on how much context is used~\cite{shannon1951}. Even with the
pessimistic $\mathcal{H} = 1$, a context of $\abs{q} = 21$ letters (the shortest one
\cref{alg:settle} tries) in a window of $\abs{S} = 2000$ letters gives
$\E[\text{spurious}] \approx 2000 \cdot 2^{-21} \approx 10^{-3}$. Spurious
matches come instead from genuine repetition, such as a phrase used twice in
nearby paragraphs, which is exactly the case the distance
\cref{eq:line-distance} resolves.

\subsection{Selecting whole words}

Let a drag select characters $[s, t)$ of the per-word text layer of
\cref{thm:per-word}. Papilot widens the selection to whole words, but only at
a boundary that cuts through a word:
\begin{equation}
  s' = \begin{cases}
    \max\{ i \le s : \text{$i = 0$ or $c_{i-1}$ is not a letter} \}, &
      c_s \text{ is a letter},\\
    s, & \text{otherwise},
  \end{cases}
  \label{eq:whole-words}
\end{equation}
and symmetrically for $t$. A drag that ends at the start of the next word, as
browsers often report, therefore keeps exactly the words it covers. The note
stores the widened quote, one rectangle per printed line, and the source lines
of \cref{alg:settle} together with an excerpt of their \LaTeX{}, so that the
agent receives formulas as source rather than as extracted text.
